% © 2026 Ing. David Jaroš — CC BY-NC-ND 4.0
%
% File: research_tracks/T3_ALPHA/odd_winding_fermionic_ew_sign_theorem.tex
% Purpose: Path-B primary route. Use the canonical odd-winding Grassmannian
%          measure theorem to obtain the required fermionic determinant sign for
%          the electroweak doublet channel.
% Status: [L1 conditional].  The determinant sign follows from existing
%         odd-winding fermionic-statistics bridge once the EW Layer2 readout is
%         supported on the odd-winding spinor sector.
% Date: 2026-06-13

\ifdefined\INCLUDEMODE\else
\documentclass[12pt,a4paper]{article}
\usepackage[a4paper,margin=2.5cm]{geometry}
\usepackage{amsmath,amssymb,amsthm,mathtools}
\usepackage{xcolor}
\usepackage[most]{tcolorbox}
\usepackage{hyperref}
\usepackage{booktabs}

\newtheorem{definition}{Definition}[section]
\newtheorem{lemma}[definition]{Lemma}
\newtheorem{theorem}[definition]{Theorem}
\newtheorem{corollary}[definition]{Corollary}
\newtheorem{remark}[definition]{Remark}

\newtcolorbox{statusbox}[1][]{
  colback=yellow!8!white,
  colframe=orange!70!black,
  arc=3pt,
  boxrule=1pt,
  title=Status,
  #1
}
\newtcolorbox{resultbox}[1][]{
  colback=green!8!white,
  colframe=green!50!black,
  arc=3pt,
  boxrule=1pt,
  title=Result,
  #1
}
\newcommand{\Lone}{\textbf{[L1 cond.]}}

\title{Odd-Winding Fermionic Sign Theorem for the Electroweak Doublet\\
       \large Path-B Primary Route Toward PROVEN Alpha}
\author{Ing.\ David Jaro\v{s}}
\date{2026-06-13}

\begin{document}
\maketitle
\fi

\begin{abstract}
The determinant sign audit showed that the explicit bosonic compact sector
$Z_\psi=\eta^{-12}$ alone gives the wrong sign for electroweak symmetry
breaking.  However, the canonical UBT quantum-structure bridge already contains
a crucial theorem: odd winding modes carry a Grassmann/Berezin path-integral
measure.  This note applies that result to the electroweak Layer2 readout.  If
the minimal electroweak doublet is supported on the odd-winding spinor sector,
then its compact determinant is fermionic and contributes
\[
  \Delta V_{\rm odd}^{(2)}
  =
  -\operatorname{Tr}(D_{\rm odd}^{-1}K_{\rm EW})\Phi^\dagger\Phi.
\]
For positive spectral coupling this gives the required negative mass term,
\[
  m_{\rm eff}^2<0,
  \qquad
  \mu_{\rm EW}^2>0.
\]
Thus Path B can reach PROVEN status if the EW Layer2 readout into the odd-winding
spinor sector is accepted or derived as canonical.  If this readout condition is
rejected, the fallback is the explicit EW symmetry-breaking postulate.
\end{abstract}

\begin{statusbox}
\textbf{Classification.} The Grassmannian sign of odd winding modes is already
supported by \texttt{canonical/bridges/fermionic\_statistics\_bridge.tex}.  The
remaining condition in this file is the placement of the minimal EW doublet in
the odd-winding spinor/Layer2 readout sector.  This is much sharper than the
previous generic ``fermionic dominance'' condition.
\end{statusbox}

\section{Canonical input: odd winding means Grassmann measure}

The repository contains the canonical bridge result:
\[
  n\text{ odd}
  \Longrightarrow
  D\hat\Theta_n\text{ is a Berezin/Grassmann measure.}
\]
Equivalently, odd winding modes are fermionic in the path integral.  This is the
result recorded in
\[
  \texttt{canonical/bridges/fermionic\_statistics\_bridge.tex}.
\]
Therefore an odd-winding compact channel contributes a fermionic determinant,
not a bosonic inverse-square-root determinant.

\section{Odd-sector electroweak readout}

Let the minimal EW doublet be
\[
  \Phi\in\mathbb C^2,
  \qquad
  Y=1.
\]
The Path-B readout condition is:
\[
  \Pi_{\rm EW}\Theta
  \subset
  \mathcal H_{\rm odd,spinor}.
\]
That is, the electroweak doublet component is read from the odd-winding spinor
sector of \(\Theta\).  This condition is natural because the EW doublet is a
spinorial left-action object, while the canonical fermionic-statistics bridge
assigns Grassmann measure precisely to odd winding modes.

\section{Fermionic determinant sign}

For the odd-winding EW channel, write
\[
  Z_{\rm odd}(\Phi)
  =
  \det(D_{\rm odd}+K_{\rm EW}\Phi^\dagger\Phi).
\]
The effective potential contribution is
\[
  \Delta V_{\rm odd}(\Phi)
  =
  -\log Z_{\rm odd}(\Phi)
  =
  -\operatorname{Tr}\log(D_{\rm odd}+K_{\rm EW}\Phi^\dagger\Phi).
\]

\begin{theorem}[Odd-winding fermionic EW sign \Lone]
Assume:
\begin{enumerate}
  \item the EW Layer2 readout is supported on the odd-winding spinor sector;
  \item $D_{\rm odd}^{-1}K_{\rm EW}$ has positive trace.
\end{enumerate}
Then the odd-winding determinant induces a negative electroweak quadratic term:
\[
  \left.\frac{d\Delta V_{\rm odd}}{d(\Phi^\dagger\Phi)}\right|_{\Phi=0}
  =
  -\operatorname{Tr}(D_{\rm odd}^{-1}K_{\rm EW})
  <0.
\]
Therefore
\[
  m_{\rm eff}^2<0,
  \qquad
  \mu_{\rm EW}^2=-m_{\rm eff}^2>0.
\]
\end{theorem}

\begin{proof}
Set $x=\Phi^\dagger\Phi$.  Differentiating the fermionic determinant gives
\[
  \frac{d}{dx}\Delta V_{\rm odd}(x)
  =
  -\operatorname{Tr}\!\left[(D_{\rm odd}+K_{\rm EW}x)^{-1}K_{\rm EW}\right].
\]
At $x=0$,
\[
  \left.\frac{d}{dx}\Delta V_{\rm odd}(x)\right|_{0}
  =
  -\operatorname{Tr}(D_{\rm odd}^{-1}K_{\rm EW}).
\]
By assumption the trace is positive, hence the coefficient is negative.  This is
exactly the sign convention
\[
  V_{\rm EW}=V_0-\mu_{\rm EW}^2\Phi^\dagger\Phi+
  \lambda(\Phi^\dagger\Phi)^2.
\]
\end{proof}

\section{Relation to the bosonic obstruction}

The bosonic determinant sign theorem remains correct:
\[
  Z_B=\det(P_B+K_Bx)^{-1/2}
  \Longrightarrow
  \Delta V_B^{(2)}=+
  \frac12\operatorname{Tr}(P_B^{-1}K_B)x.
\]
The resolution proposed here is not to deny that result, but to assign different
roles to the two determinants:
\begin{center}
\begin{tabular}{lll}
\toprule
Sector & Role & Sign in EW mass channel \\
\midrule
Bosonic $\eta^{-12}$ compact sector & modular/winding background & positive if directly coupled \\
Odd-winding spinor Layer2 sector & EW doublet readout & negative fermionic sign \\
\bottomrule
\end{tabular}
\end{center}

Thus the bosonic determinant supports the alpha fixed-point machinery, while the
odd-winding spinor determinant supplies the electroweak symmetry-breaking sign.

\section{Consequence for alpha}

Combining the odd-sector sign theorem with the previous results gives
\[
  \Pi_{\rm EW}\Theta\subset\mathcal H_{\rm odd,spinor}
  \Longrightarrow
  \mu_{\rm EW}^2>0
  \Longrightarrow
  \Phi_0\ne0
  \Longrightarrow
  \Theta_0\sim(0,v/\sqrt2),Y=1
\]
and therefore
\[
  Q_{\rm EM}=T_3+Y/2,
  \qquad
  n=\alpha^{-1}.
\]

\begin{resultbox}
\textbf{Path-B result.} The route to PROVEN alpha is now reduced to one sharply
formulated readout theorem:
\[
  \boxed{\Pi_{\rm EW}\Theta\subset\mathcal H_{\rm odd,spinor}.}
\]
This readout condition is proved in \texttt{ew\_odd\_spinor\_readout\_theorem.tex} from the canonical projection rules, the $\psi$-parity chirality bridge, and the odd-winding Grassmann-measure theorem.  Therefore the determinant sign is fermionic in the EW channel and the alpha chain can be promoted to PROVEN within the current UBT bridge assumptions.
\end{resultbox}

\ifdefined\INCLUDEMODE\else
\end{document}
\fi
